% =============================================================================
%  Relatório de Iniciativa — InovaDados / Poli Júnior
%
%  NÃO EDITE A CÓPIA GERADA. Edite este template em templates/latex/relatorio.tex
%  e regenere com a skill `documentar-iniciativa`.
%
%  Compilar (duas vezes, para o sumário):
%      pdflatex relatorio.tex && pdflatex relatorio.tex
% =============================================================================
\documentclass[12pt,a4paper]{article}

\usepackage[utf8]{inputenc}
\usepackage[T1]{fontenc}
\usepackage[brazil]{babel}
\usepackage{lmodern}
\usepackage[a4paper,margin=2.5cm,headheight=15pt]{geometry}
\usepackage{graphicx}
\usepackage{eso-pic}
\usepackage{tikz}
\usepackage{fancyhdr}
\usepackage{titlesec}
\usepackage{enumitem}
\usepackage{xcolor}
\usepackage{longtable}
\usepackage{booktabs}
\usepackage{microtype}
\usepackage[hidelinks]{hyperref}

% --- identidade -------------------------------------------------------------
\definecolor{pjazul}{HTML}{0B2C5B}
\definecolor{pjcinza}{HTML}{5A6472}

% --- variáveis preenchidas pelo gerador -------------------------------------
\newcommand{\nomeIniciativa}{Modelo de Precificação NDados}
\newcommand{\subtituloRelatorio}{Inovação 2026.1}
\newcommand{\marcaDagua}{assets/logo-polijunior.png}   % copiado para assets/ pelo gerador

% --- marca d'água em todas as páginas ---------------------------------------
% Opacidade via TikZ: funciona em pdflatex, xelatex e lualatex sem o pacote
% `transparent`, que depende de driver.
\AddToShipoutPictureBG{%
  \begin{tikzpicture}[remember picture,overlay]
    \node[opacity=0.06] at (current page.center)
      {\includegraphics[width=0.65\paperwidth]{\marcaDagua}};
  \end{tikzpicture}%
}

% --- cabeçalho e rodapé -----------------------------------------------------
\pagestyle{fancy}
\fancyhf{}
\renewcommand{\headrulewidth}{0.4pt}
\fancyhead[L]{\small\color{pjcinza}\nomeIniciativa}
\fancyhead[R]{\small\color{pjcinza}\subtituloRelatorio}
\fancyfoot[C]{\small\color{pjcinza}\thepage}

% --- títulos ----------------------------------------------------------------
\titleformat{\section}{\Large\bfseries\color{pjazul}}{\thesection}{1em}{}
\titleformat{\subsection}{\large\bfseries\color{pjazul}}{\thesubsection}{1em}{}
\titleformat{\subsubsection}{\normalsize\bfseries\color{pjcinza}}{}{0em}{}
\setcounter{secnumdepth}{2}   % subsubsection sem número, como no titleformat acima
\setcounter{tocdepth}{2}      % sumário até subseção
\setlist[itemize]{leftmargin=1.4em,itemsep=0.25em,topsep=0.3em}

% --- bloco de metadados de ata / framework ----------------------------------
\newcommand{\metadados}[1]{%
  \begingroup\small\color{pjcinza}%
  \begin{list}{}{\setlength{\leftmargin}{1em}\setlength{\itemsep}{0pt}}
  \item[] #1
  \end{list}%
  \endgroup\vspace{0.6em}%
}

\begin{document}

% ---------- PÁGINA 1: capa --------------------------------------------------
\thispagestyle{empty}
\vspace*{\fill}
\begin{center}
  {\fontsize{30}{36}\selectfont\bfseries\color{pjazul}\nomeIniciativa\par}
  \vspace{1.2em}
  {\Large\color{pjcinza}\subtituloRelatorio\par}
\end{center}
\vspace*{\fill}
\clearpage

% ---------- PÁGINA 2: sumário ----------------------------------------------
\thispagestyle{empty}
\renewcommand{\contentsname}{Sumário}
\tableofcontents
\clearpage

% ---------- PÁGINA 3 em diante: conteúdo gerado -----------------------------
\setcounter{page}{1}
% Gerado por .claude/skills/documentar-iniciativa — NÃO EDITE À MÃO.
% A fonte é o markdown do card; regenere após alterá-lo.

\section{Resumo Executivo}

O NDados precificava por intuição do gerente, com dispersão de até 3x entre projetos de escopo equivalente. O card entregou um modelo de três fatores — complexidade, esforço e risco — com pesos calibrados por esfera de atuação e piso de margem fixo. O backtest sobre 38 propostas históricas reproduziu o preço praticado dentro de 15\% em 29 casos; os 9 desvios concentraram-se em projetos com escopo aberto. Aplicado a três propostas reais, o modelo subiu o preço médio em 18\% sem perda de conversão. Fica pendente calibrar a esfera de engenharia de dados, sustentada por poucos casos.

% Gerado por .claude/skills/documentar-iniciativa — NÃO EDITE À MÃO.
% A fonte é o markdown do card; regenere após alterá-lo.

\clearpage
\section{Introdução}

\subsection{Contexto}

Este documento consolida os artefatos produzidos pela iniciativa \textbf{Modelo de Precificação NDados} durante o ciclo 2026.1 do núcleo NDados. Reúne o resumo executivo, as atas de benchmark realizadas e os frameworks de discovery derivados delas.

A governança da iniciativa (status, prioridade, checklist e alocação) permanece no Notion. O que está aqui é o registro de campo e a conclusão derivada dele — o material destinado a ser lido daqui a um ano.

\subsection{Por que esta iniciativa existe}

Padronizar esferar de atuação da empresa.

\subsection{Objetivo}

Conceber um modelo matemático de precificação para padronizar o preço de projetos do núcleo.

\subsection{Métricas de sucesso}

faturamento total, taxa de conversão.

\subsection{Como ler este documento}

As atas reproduzem o registro bruto das conversas, separando o que foi dito (seção \emph{Perguntas}) da interpretação da equipe (seção \emph{Síntese}). Os frameworks que vêm depois são análises derivadas: cada um declara, no campo \texttt{origem}, os IDs das atas que o embasam. Nenhuma conclusão apresentada nos frameworks existe sem uma evidência rastreável nas atas.

% Gerado por .claude/skills/documentar-iniciativa — NÃO EDITE À MÃO.
% A fonte é o markdown do card; regenere após alterá-lo.

\clearpage
\section{Ata bm-001 — Rafael Toledo — Núcleo de Dados, EJ Sigma (fictícia)}

\metadados{\textbf{entrevistado:} Rafael Toledo — Núcleo de Dados, EJ Sigma (fictícia) \\ \textbf{entrevistador:} Helena Braga \\ \textbf{data:} 2026-03-04}


\subsection{Perguntas}

\subsubsection{Como vocês chegam ao preço de um projeto de dados hoje?}

Partimos de uma estimativa de horas por perfil — analista, engenheiro, gerente — e multiplicamos por uma taxa interna. A taxa é revisada uma vez por ano, no planejamento. Em cima disso entra um ajuste que o gerente comercial faz “no olho”, que ele chama de fator de dor: o quanto o cliente parece sofrer com o problema.

\subsubsection{Esse ajuste no olho é registrado em algum lugar?}

Não. É a maior fragilidade do processo. Quando o gerente sai, o critério sai com ele. Já aconteceu de dois projetos quase idênticos saírem com 40\% de diferença de preço em semestres seguidos, e ninguém conseguiu reconstruir por quê.

\subsubsection{Vocês diferenciam o preço por tipo de projeto?}

Diferenciamos por duração, não por natureza. Na prática, um dashboard e um modelo preditivo de mesma duração saem pelo mesmo preço, o que é claramente errado — o preditivo consome muito mais sênior e tem risco de não convergir.

\subsubsection{O que vocês já tentaram e não funcionou?}

Tentamos precificação por valor gerado, pedindo ao cliente a estimativa de ganho. Não funcionou: o cliente não tem esse número, e quando tem não compartilha. Voltamos para custo mais margem em dois meses.

\subsubsection{Como vocês tratam projeto de escopo aberto?}

Mal. Fechamos preço fixo e absorvemos o estouro. Foi a maior fonte de prejuízo do último ciclo.

\subsection{Síntese}

\subsubsection{Objeções}

\begin{itemize}
  \item Precificação por valor gerado foi testada e abandonada: depende de um número que o cliente não tem ou não compartilha.
  \item Resistência interna a “engessar” o julgamento comercial em fórmula.
\end{itemize}

\subsubsection{Surpresas}

\begin{itemize}
  \item A taxa interna é revisada só uma vez por ano, mesmo com o custo do time mudando a cada virada de time.
  \item Duração é o único eixo de diferenciação; natureza do projeto não entra.
\end{itemize}

\subsubsection{Oportunidades}

\begin{itemize}
  \item O fator de dor é tácito e não registrado — transformá-lo em parâmetro explícito resolve o problema de sucessão que eles descreveram.
  \item Separar preço por natureza do projeto é ganho imediato e barato.
\end{itemize}

\subsubsection{Ameaças}

\begin{itemize}
  \item Escopo aberto com preço fixo é a fonte de prejuízo deles; qualquer modelo nosso que ignore isso repete o erro.
  \item Modelo percebido como camisa de força é abandonado pelo comercial em poucos meses.
\end{itemize}


\clearpage
\section{Ata bm-002 — Marina Cerqueira — Consultoria Delta (fictícia)}

\metadados{\textbf{entrevistado:} Marina Cerqueira — Consultoria Delta (fictícia) \\ \textbf{entrevistador:} Helena Braga \\ \textbf{data:} 2026-03-11}


\subsection{Perguntas}

\subsubsection{Vocês usam um modelo formal de precificação?}

Usamos, há três anos. É uma planilha com três fatores: complexidade técnica, volume de esforço e risco. Cada fator vira uma nota de 1 a 5, e o produto das três entra numa tabela de faixas. O comercial não define preço, define as notas.

\subsubsection{Como vocês evitam que o comercial infle as notas?}

A nota de risco é dada pelo time técnico, não pelo comercial. Esse foi o ajuste que fez o modelo parar de ser burlado. Antes, quem vendia dava as três notas e o resultado sempre caía na faixa que ele já tinha na cabeça.

\subsubsection{O modelo acerta?}

Acerta o suficiente. Comparamos o preço do modelo com o preço fechado e ficamos dentro de 10\% na maioria dos casos. Onde ele erra feio é em projeto de descoberta, onde o escopo muda no meio — aí nenhuma fórmula salva.

\subsubsection{Quanto tempo levou para o time confiar no modelo?}

Cerca de dois ciclos. O que destravou foi rodar o modelo em paralelo, sem poder de veto, durante um semestre inteiro. As pessoas viram que ele não estava longe do que elas já fariam, e aí pararam de brigar.

\subsubsection{Se fossem começar hoje, o que fariam diferente?}

Definiríamos o piso de margem antes da fórmula, não depois. Passamos um ano com um modelo que conseguia sugerir preço abaixo do nosso custo em certas combinações de nota, e só descobrimos porque alguém reparou.

\subsubsection{Qual o custo de manter isso?}

Baixo, mas não zero. Uma revisão de pesos por ciclo, olhando os desvios acumulados. Quem larga a revisão volta ao chute em um ano.

\subsection{Síntese}

\subsubsection{Objeções}

\begin{itemize}
  \item Modelo formal só é aceito depois de rodar em paralelo, sem poder de veto, por um ciclo inteiro — adoção direta gera resistência.
  \item Nenhum modelo resolve projeto de descoberta com escopo móvel.
\end{itemize}

\subsubsection{Surpresas}

\begin{itemize}
  \item A separação de quem dá cada nota importa mais que a fórmula: risco avaliado pelo time técnico foi o que impediu a manipulação.
  \item Piso de margem definido depois da fórmula deixou um ano de propostas abaixo do custo passarem despercebidas.
\end{itemize}

\subsubsection{Oportunidades}

\begin{itemize}
  \item Estrutura de três fatores com nota de 1 a 5 é simples o bastante para o nosso comercial operar sem treinamento longo.
  \item Revisão de pesos por ciclo cabe no nosso calendário de virada de time.
\end{itemize}

\subsubsection{Ameaças}

\begin{itemize}
  \item Sem revisão periódica, o modelo apodrece e o time volta ao chute em um ano.
  \item Se o mesmo papel der as três notas, o modelo vira justificativa do preço que já se queria praticar.
\end{itemize}


% Gerado por .claude/skills/documentar-iniciativa — NÃO EDITE À MÃO.
% A fonte é o markdown do card; regenere após alterá-lo.

\clearpage
\section{Matriz CSD (fw-001)}

\metadados{\textbf{tipo:} csd \\ \textbf{origem:} bm-001, bm-002}


\subsubsection{Certezas}

\begin{itemize}
  \item Precificação baseada em julgamento não registrado não sobrevive à troca de gerente. Os dois entrevistados relataram o mesmo problema de sucessão (\texttt{bm-001}, \texttt{bm-002}).
  \item Duração isolada não explica preço: projetos de mesma duração e naturezas diferentes consomem perfis de senioridade muito distintos (\texttt{bm-001}).
  \item Um modelo de três fatores com nota de 1 a 5 é operável por um comercial sem treinamento longo — está em produção há três anos na Delta (\texttt{bm-002}).
  \item Escopo aberto com preço fixo gera prejuízo recorrente (\texttt{bm-001}).
  \item O piso de margem precisa ser definido antes da fórmula, não depois (\texttt{bm-002}).
\end{itemize}

\subsubsection{Suposições}

\begin{itemize}
  \item A base histórica do NDados tem volume suficiente para calibrar os pesos das quatro esferas. A checar no backtest — desconfiamos que engenharia de dados tenha poucos casos.
  \item Nosso comercial aceitará operar por notas se o modelo rodar em paralelo por um ciclo, como funcionou na Delta. Nada garante que a cultura seja transferível.
  \item A dispersão de preço que observamos internamente vem da mesma causa relatada no benchmark, e não de diferenças reais de escopo mal registradas no CRM.
  \item Separar quem dá a nota de risco de quem dá as outras duas é suficiente para evitar a manipulação no nosso contexto, com um time menor que o da Delta.
\end{itemize}

\subsubsection{Dúvidas}

\begin{itemize}
  \item Como precificar projeto de descoberta, em que nenhum dos dois entrevistados tem resposta e ambos relatam prejuízo?
  \item Qual a periodicidade certa de revisão dos pesos para um núcleo que troca de time duas vezes por ano, contra o ciclo anual da Delta?
  \item O fator de risco deve multiplicar o preço ou definir um percentual de contingência separado, visível na proposta?
  \item Quem é o dono do modelo depois que este card fecha?
\end{itemize}


\clearpage
\section{Análise SWOT (fw-002)}

\metadados{\textbf{tipo:} swot \\ \textbf{origem:} bm-001, bm-002}


\subsubsection{Strengths}

\begin{itemize}
  \item Base histórica de 38 propostas já normalizada por horas-equipe, o que nenhum dos dois núcleos entrevistados tinha ao começar (\texttt{bm-001}, \texttt{bm-002}).
  \item Núcleo pequeno o bastante para que a separação de papéis na atribuição de notas seja aplicável sem burocracia.
  \item Time com capacidade técnica para manter o backtest como código, não como planilha.
\end{itemize}

\subsubsection{Weaknesses}

\begin{itemize}
  \item A esfera de engenharia de dados está sustentada por poucos casos históricos; os pesos dela são os menos confiáveis do modelo.
  \item Troca de time duas vezes por ano ameaça a continuidade da revisão de pesos, que a Delta faz com time estável (\texttt{bm-002}).
  \item Nenhum dono formal designado para o modelo depois do fechamento deste card.
\end{itemize}

\subsubsection{Opportunities}

\begin{itemize}
  \item Transformar o fator de dor tácito em parâmetro explícito resolve o problema de sucessão relatado no benchmark e é diferencial frente a EJs comparáveis (\texttt{bm-001}).
  \item Rodar o modelo em paralelo por um ciclo, sem poder de veto, replica o caminho de adoção que funcionou na Delta e custa pouco (\texttt{bm-002}).
  \item Preço diferenciado por natureza do projeto abre espaço para subir a margem em preditivo, hoje precificado como dashboard.
\end{itemize}

\subsubsection{Threats}

\begin{itemize}
  \item Modelo percebido como camisa de força é abandonado pelo comercial em poucos meses (\texttt{bm-001}).
  \item Sem revisão periódica dos pesos, o modelo apodrece e o time volta ao chute em um ano (\texttt{bm-002}).
  \item Projeto de escopo aberto continua sem tratamento: aplicar o modelo nele dá falsa sensação de rigor sobre um número que não se sustenta.
\end{itemize}



\end{document}
